\documentclass[10pt]{article}
\usepackage[T1]{fontenc}
\usepackage[utf8]{inputenc}
\usepackage[french]{babel}
\usepackage{amsmath,amssymb,mathrsfs,mathpazo}
\usepackage{geometry}
\geometry{a4paper,margin=2cm}
\usepackage{graphicx,array,tabularx,tabularray,colortbl,multirow,multicol,booktabs,enumitem,pifont,ulem,fancybox,diagbox,cancel}
\usepackage{pstricks,pstricks-add,pst-plot,pst-tree,pst-node,pst-text,pst-func,pst-math,pst-bspline,pst-3dplot,pst-eucl,pst-solides3d}
\let\pscalculate\undefined
\usepackage{pst-bezier}
\usepackage{tikz,pgfplots,tkz-tab}
\usetikzlibrary{arrows,arrows.meta,patterns,calc,positioning,trees,decorations.markings,shapes.geometric,shapes.misc,shapes.arrows,intersections,decorations.pathmorphing,decorations.pathreplacing,backgrounds,shadows,babel}
\pgfplotsset{compat=1.18}
\usepgfplotslibrary{fillbetween}
\usepackage[np]{numprint}
\usepackage{listings,makecell,pas-tableur,twoopt}
\let\barre\undefined
\usepackage[pstricks]{tablvar}
\definecolor{darkgreen}{RGB}{0,100,0}
\providecommand{\red}{\color{red}}
\providecommand{\blue}{\color{blue}}
\providecommand{\green}{\color{green}}
\providecommand{\black}{\color{black}}
\providecommand{\white}{\color{white}}
\providecommand{\gray}{\color{gray}}
\providecommand{\cyan}{\color{cyan}}
\providecommand{\magenta}{\color{magenta}}
\providecommand{\yellow}{\color{yellow}}
\providecommand{\violet}{\color{violet}}
\providecommand{\orange}{\color{orange}}
\providecommand{\euro}{\texteuro}
\providecommand{\EUR}{\texteuro}
\providecommand{\ieme}{\textsuperscript{e}}
\providecommand{\degres}{\ensuremath{^\circ}}
\providecommand{\degree}{\ensuremath{^\circ}}
\providecommand{\vv}[1]{\overrightarrow{#1}}
\providecommand{\Coord}[1]{#1}
\providecommand{\decosix}{}
\providecommand{\starredbullet}{\ensuremath{\bullet}}
\pagestyle{empty}
\setlength{\parindent}{0pt}
\usepackage{adjustbox}\DeclareUnicodeCharacter{2236}{:}\DeclareUnicodeCharacter{2190}{\ensuremath{\leftarrow}}
\providecommand{\R}{\mathbb{R}}
\providecommand{\N}{\mathbb{N}}
\providecommand{\Z}{\mathbb{Z}}
\providecommand{\Q}{\mathbb{Q}}
\providecommand{\D}{\mathbb{D}}
\providecommand{\C}{\mathbb{C}}
\providecommand{\e}{{\mathrm{e}}}
\providecommand{\vect}[1]{\overrightarrow{#1}}
\providecommand{\vectt}[1]{\overrightarrow{\mathrm{#1}}}
\providecommand{\vvt}[1]{\overrightarrow{\mathrm{#1}}}
\providecommand{\np}[2][]{{#2~#1}}
\providecommand{\vv}[1]{{\overrightarrow{#1}}}
\renewcommand{\d}{{\mathrm{d}}}
\providecommand{\degres}{{^\circ}}
\providecommand{\degre}{{^\circ}}
\providecommand{\blue}{\color{blue}}
\providecommand{\red}{\color{red}}
\providecommand{\enskip}{\quad}
\providecommand{\llbracket}{\lbrack\!\lbrack}
\providecommand{\rrbracket}{\rbrack\!\rbrack}
\providecommand{\euro}{{\text{€}}}
\providecommand{\no}{{\text{n°}}}
\providecommand{\interval}[2]{{[#1;#2]}}
\providecommand{\intervFO}[2]{{[#1;#2[}}
\providecommand{\dotfill}{\ldots}

\providecommand{\ds}{\displaystyle}
\providecommand{\textding}[1]{#1}
\providecommand{\Oij}{\(\mathrm{O};\vec{\imath},\vec{\jmath}\)}
\providecommand{\Oijk}{\(\mathrm{O};\vec{\imath},\vec{\jmath},\vec{k}\)}
\providecommand{\Ouv}{\(\mathrm{O};\vec{u},\vec{v}\)}
\providecommand{\intervFF}[2]{\ensuremath{[#1;#2]}}
\renewcommand{\euro}{\ifmmode\text{\texteuro}\else\texteuro\fi}
\DeclareUnicodeCharacter{00B0}{\ifmmode^{\circ}\else\textdegree\fi}
\DeclareUnicodeCharacter{20AC}{\ifmmode\text{\texteuro}\else\texteuro\fi}
% Énoncés : Éducation nationale. Transcription et corrigés : APMEP, Denis Vergès et contributeurs. Découpage et classement : CoopMaths, Sébastien Lozano, Éric Elter et contributeurs. Adaptation de la présentation : Maths Explorer.
% CC BY-SA : https://coopmaths.fr/www/about/
% https://coopmaths.fr/alea/static/bac/2026/tex/bac_2026_05_sujet1_ameriquenord_3_cor.tex
\begin{document}
\noindent\textit{APMEP / CoopMaths. CC BY-SA. Présentation adaptée par Maths Explorer.}\par\medskip
\textbf{Partie A}

\begin{enumerate}
\item 

Dans le repère $\left(\mathrm{O}; \vectt{OI}, \vectt{OJ}, \vectt{OK}\right)$ on a :\quad A$(-1~;~-1~;~0)$ \quad et \quad D$(1~;~-1~;~0)$.

\item Étant données les coordonnées de S et C, on a donc :

$\vectt{SC}\begin{pmatrix}	x_\mathrm{C}-x_\mathrm{S} \\ y_\mathrm{C}-y_\mathrm{S}\\ z_\mathrm{C}-z_\mathrm{S} \end{pmatrix}$ \quad soit \quad $\vectt{SC}\begin{pmatrix}	1-0 \\ 1-0\\ 0-2 \end{pmatrix}$ \quad finalement :\quad $\vectt{SC}\begin{pmatrix}	1 \\1\\-2 \end{pmatrix}$
\quad et de même :\quad $\vectt{SB}\begin{pmatrix}	-1 \\1\\-2 \end{pmatrix}$

On est dans un repère orthonormé donc :

$\vectt{SC}\cdot\vectt{SB}= 1 \times (-1) + 1 \times 1 + (-2) \times (-2) = -1 +1+4=4$.

\item On a SC$=\sqrt{ 1^2+1^2+(-2)^2} = \sqrt{6}$ \quad et \quad SB$=\sqrt{ (-1)^2+1^2+(-2)^2} = \sqrt{6}$

$	\begin{aligned} \vectt{SC}\cdot\vectt{SB} = \mathrm{SC} \times \mathrm{SB} \times \cos\left(\widehat{\mathrm{BSC}}\right) & \iff 4 = \sqrt{6} \times \sqrt{6} \times \cos\left(\widehat{\mathrm{BSC}}\right) \\
			& \iff \dfrac{4}{6} = \cos\left(\widehat{\mathrm{BSC}}\right) \\
				& \iff \dfrac{2}{3} = \cos\left(\widehat{\mathrm{BSC}}\right) \\
			\end{aligned} $

	D'où \quad $\widehat{\mathrm{BSC}} = \arccos\left(\dfrac{2}{3}\right) \approx 48,1897$(\ensuremath{^\circ}) \quad soit \quad $\widehat{\mathrm{BSC}} \approx 48,2 $(\ensuremath{^\circ}),\quad arrondi au dixième de degré près.
	\end{enumerate}

	\textbf{Partie B}
	
	\begin{enumerate}
	\item \begin{enumerate}
		\item Les vecteurs $\vectt{SB}$ et $\vectt{SC}$ ne sont pas colinéaires.
		(On peut le justifier par étude des coordonnées des vecteurs, ou parce que l'angle $\widehat{\mathrm{BSC}}$ n'est ni plat ni nul, ou parce que le plan étant nommé (SBC), cela implique que les points S, B et C ne sont pas alignés).

		De plus : \quad $\vect{n} \cdot \vectt{SB} = 0 \times (-1) + 2 \times 1 +1 \times (-2) = 0 +2-2=0$

		et \quad $\vect{n} \cdot \vectt{SC} = 0 \times (1) + 2 \times 1 +1 \times (-2) = 0 +2-2=0$

		Le vecteur $\vect{n}$ est donc orthogonal à deux vecteurs non colinéaires du plan (SBC) donc le vecteur $\vect{n}$ est normal au plan (SBC).

		\item Le vecteur $\vect{n}$ étant normal au plan (SBC), une équation cartésienne de ce plan est de la forme :\quad $0x+2y+z+d=0$ \quad avec $d \in \R$.

		$S(0~;~0~;~2) \in (\mathrm{SBC})$ \quad donc \quad $0+2\times 2+d=0 \iff d=-2$.

		Donc une équation cartésienne du plan (SBC) est \quad $2 y+z-2=0$.

\end{enumerate}
\item
	\begin{enumerate}
			\item H est le projeté orthogonal du point O sur le plan (SBC) donc la droite (OH) est orthogonale au plan (SBC). Un vecteur normal au plan (SBC) sera donc un vecteur directeur de la droite (OH).

Un vecteur directeur de la droite (OH) est donc le vecteur $\vect{n}\begin{pmatrix} 0\\ 2\\ 1\\ \end{pmatrix}$.

De plus le point O, de coordonnées $(0~;~0~;~0)$ appartient à la droite (OH) donc une représentation de la droite (OH) est :\quad
$\begin{cases}
x =x_{\text{O}} + x_{\vect{n}}t\\
y =y_{\text{O}} + y_{\vect{n}}t\\
z =z_{\text{O}} + z_{\vect{n}}t\\
	\end{cases}\text{ avec } t \in \R$ \quad soit \quad
	$\begin{cases}
x =0\\
y =2 t\\
z =t\\
\end{cases}\text{ avec } t~\in~\R.$
		\item H est le projeté orthogonal de O sur (SBC), alors $\mathrm{H} = (\mathrm{OH}) \cap \mathrm{(SBC)}$.

Donc H est le point de (OH) dont le paramètre fait que les coordonnées paramétriques vérifient l'équation de (SBC).

On résout donc : 
\quad $\begin{aligned}[t]
	4t+t-2= 0 & \iff 5t = 2\\
	& \iff t = \dfrac{2}{5}
	\end{aligned}$

H est donc le point de paramètre $t = \dfrac{2}{5}$ sur (OH),

donc ses coordonnées sont : \quad $\mathrm{H} \left(0 ~;~ \dfrac{4}{5} ~;~ \dfrac{2}{5}\right)$.

\item La distance OH vaut donc : \quad $\begin{aligned}[t]
	\mathrm{OH} &= \sqrt{(0-0)^2 + \left(\dfrac{4}{5}-0 \right)^2 + \left(\dfrac{2}{5} -0 \right)^2}\\
			& =\sqrt{\dfrac{16}{25} + \dfrac{4}{25}}=\sqrt{\dfrac{20}{25}}=\sqrt{\dfrac{4 \times 5}{5^2}}\\
			&=\dfrac{2\sqrt{5}}{5}
		\end{aligned}$

	La distance du point O au plan (SBC), c'est-à-dire la distance de O à son projeté orthogonal sur (SBC), est donc bien égale à $\dfrac{2 \sqrt{5}}{5}$~(cm).
		\end{enumerate}
	\end{enumerate}

\textbf{Partie C}

\begin{enumerate}
\item
	\begin{enumerate}
			\item La base de la pyramide SABCD est le carré ABCD, de côté 2~cm et sa hauteur associé est le segment [SO], de longueur 2~cm.

Donc $V= \dfrac{1}{3} \times 2^2\times 2 = \dfrac{8}{3}$~(cm\up{3}.

Son volume est donc égal à \quad $\dfrac{8}{3}$~(cm\up{3}).

			\item La base de la pyramide OCBS est le triangle OCB et sa hauteur associé est le segment [SO].

L'aire du triangle OCB est égale au quart de l'aire du carré ABCD car O est le centre du carré ABCD.

La hauteur étant la même, le volume de la pyramide OCBS est égal au quart du volume de la pyramide SABCD.

$V'= \dfrac{1}{4} V = \dfrac{1}{4} \times \dfrac{8}{3} = \dfrac{2}{3}$.

Donc le volume de la pyramide OCBS est bien égal à $\dfrac{2}{3}$~(cm\up{3}).
	\end{enumerate}

		\item Dans le triangle SBC, on a : \quad $\mathrm{SB} = \mathrm{SC} = \sqrt{6}$ \quad donc SBC est un triangle isocèle en S.

J étant le milieu du segment [BC], le segment [SJ] est la hauteur issue de S dans le triangle SBC.

Le triangle JSB est rectangle en J donc, d'après le théorème de Pythagore,

$\begin{aligned} \mathrm{SB}^2 = \mathrm{SJ}^2+ \mathrm{BJ}^2 & \iff
\mathrm{SJ}^2 = \mathrm{SB}^2- \mathrm{BJ}^2 \\
	& \iff \mathrm{SJ}^2 = \left(\sqrt{6}\right)^2- 1^2= 6-1=5\\
	& \iff	\mathrm{SJ}=\sqrt{5} \quad \text{ car une longueur est positive}
\end{aligned}$

Soit $A$ l'aire du triangle SBC.

On a : \quad $A= \dfrac{\text{longueur de la base} \times \text{hauteur}}{2} = \dfrac{2 \times \sqrt{5}}{2}= \sqrt{5}$

L'aire du triangle SBC est donc $\sqrt{5}$~(cm\up{2}).

		\item Calculons le volume de la pyramide OCBS en prenant le triangle SBC comme base. La hauteur associée sera donc le segment [OH] car H est le projeté orthogonal de O sur le plan (SBC).

On a donc : \quad $ V_{\mathrm{OCBS}} = \dfrac{1}{3} \times \sqrt{5} \times \mathrm{OH}$

Or, d'après la question \textbf{1. b.} de la \textbf{Partie C}, \quad$ V_{\mathrm{OCBS}} = \dfrac{2}{3}$~(cm\up{3}).

Donc \quad $\dfrac{2}{3}= \dfrac{\sqrt{5}}{3} \mathrm{OH} \iff \mathrm{OH} = \dfrac{2}{3} \times \dfrac{3}{\sqrt{5}} = \dfrac{2}{\sqrt{5}} = \dfrac{2 \sqrt{5}}{\sqrt{5}\times \sqrt{5}}x= \dfrac{2\sqrt{5}}{5}$.

Donc la distance du point O au plan (SBC), c'est-à-dire la distance de O à son projeté orthogonal sur (SBC), est bien égale à $\dfrac{2 \sqrt{5}}{5}$~(cm).

Cela confirme bien le résultat de la \textbf{Partie B}.
	\end{enumerate}

\end{document}
